% Method Module A: Problem Formulation
% This snippet is designed to be pasted into the Method section after user confirmation.

\subsection{Problem Formulation}

We treat cross-cultural recommendation as a structured downstream learning problem built on top of frozen backbone embeddings rather than as direct nearest-neighbor retrieval in the raw embedding space.
The training catalog is represented as
%
% code map: dcas/data/npz_tracks.py:10-15, 33-53
\begin{align}
\mathcal{D}_{\mathrm{trk}}
=
\left\{
\left(
t_i,
c_i,
\mathbf{e}_i,
y_i^{\mathrm{aff}},
s_i
\right)
\right\}_{i=1}^{N},
\end{align}
%
where $t_i$ is the track identifier, $c_i \in \mathcal{C}$ is the culture label, $\mathbf{e}_i \in \mathbb{R}^{D}$ is the frozen backbone embedding, $y_i^{\mathrm{aff}}$ is an optional affect label, and $s_i$ is an optional source-dataset label. In the implementation, these fields correspond directly to `track_id`, `culture`, `embedding`, `affect_label`, and `source_dataset` in the `Tracks` container.

During optimization, the loader converts the catalog into mini-batches
%
% code map: dcas/utils.py:24-30; dcas/data/torch_dataset.py:57-82; dcas/data/batch.py:10-33
\begin{align}
\mathcal{B}
=
\left(
\mathbf{X},
\mathbf{c},
\mathbf{r},
\mathbf{y}^{\mathrm{aff}},
\mathbf{s}
\right),
\end{align}
%
where $\mathbf{X} \in \mathbb{R}^{B \times D}$ is the batch embedding matrix, $\mathbf{c} \in \{1,\dots,|\mathcal{C}|\}^{B}$ is the culture index vector, $\mathbf{r}$ stores track indices, and $\mathbf{y}^{\mathrm{aff}}$ and $\mathbf{s}$ are optional affect and source labels. In code, these are the tensors `batch.x`, `batch.culture`, `batch.track_index`, `batch.affect_label`, and `batch.source_label`.

The downstream model learns a factorized latent map
%
% code map: dcas/models/dcas_vae.py:24-47, 49-145
\begin{align}
f_{\theta} : \mathbf{e}_i \mapsto
\left(
\mathbf{z}^{c}_{i},
\mathbf{z}^{s}_{i},
\mathbf{z}^{a}_{i}
\right),
\end{align}
%
with $\mathbf{z}^{c}_{i} \in \mathbb{R}^{d_c}$, $\mathbf{z}^{s}_{i} \in \mathbb{R}^{d_s}$, and $\mathbf{z}^{a}_{i} \in \mathbb{R}^{d_a}$ denoting content, style, and affective-functional latent variables, respectively.
For the current V4 main operating point, the default dimensions are $(d_c,d_s,d_a)=(32,32,16)$, matching the code parameters `zc_dim`, `zs_dim`, and `za_dim` in `DCASConfig`.

The recommendation problem is defined over a user history
%
% code map: dcas/pipelines.py:47-110, 229-276
\begin{align}
H_u = \{(i, w_{ui})\}_{i \in \mathcal{I}_u},
\end{align}
%
where $\mathcal{I}_u$ is the set of tracks observed for user $u$ and $w_{ui}$ is the interaction weight read from the synthesized interaction log. Given a target culture $c^\star \in \mathcal{C}$, the candidate set is
\begin{align}
\mathcal{X}(c^\star) = \{ i \mid c_i = c^\star \},
\end{align}
and the goal is to return a ranked list from $\mathcal{X}(c^\star)$ that remains behaviorally relevant while improving cross-cultural discovery and calibrated minority exposure.

The full training objective is stage-dependent. We first define the base representation objective
%
% code map: dcas/models/dcas_vae.py:165-203
\begin{align}
\mathcal{L}_{\mathrm{base}}
&=
\mathcal{L}_{\mathrm{recon}}
 + \beta_{\mathrm{KL}} \mathcal{L}_{\mathrm{KL}}
 + \lambda_{\mathrm{dom}}\, g_{\mathrm{dom}}(t)\, \mathcal{L}_{\mathrm{dom}}
 + \lambda_{\mathrm{con}}\, g_{\mathrm{con}}(t)\, \mathcal{L}_{\mathrm{con}} \nonumber\\
&\quad
 + \lambda_{\mathrm{cov}}\, g_{\mathrm{cov}}(t)\, \mathcal{L}_{\mathrm{cov}}
 + \lambda_{\mathrm{tc}}\, g_{\mathrm{tc}}(t)\, \mathcal{L}_{\mathrm{tc}}
 + \lambda_{\mathrm{hsic}}\, g_{\mathrm{hsic}}(t)\, \mathcal{L}_{\mathrm{hsic}} \nonumber\\
&\quad
 + \lambda_{\mathrm{aff}}\, g_{\mathrm{aff}}(t)\, \mathcal{L}_{\mathrm{aff}}
 + \lambda_{\mathrm{src}}\, g_{\mathrm{src}}(t)\, \mathcal{L}_{\mathrm{src}},
\end{align}
%
where $\mathcal{L}_{\mathrm{recon}}$ is the reconstruction loss, $\mathcal{L}_{\mathrm{KL}}$ is the sum of three standard-normal KL terms, $\mathcal{L}_{\mathrm{dom}}$ is the culture-adversarial loss on $\mathbf{z}^a$, $\mathcal{L}_{\mathrm{con}}$ is the augmentation-based InfoNCE loss on $\mathbf{z}^c$, $\mathcal{L}_{\mathrm{cov}}$ penalizes off-diagonal covariance, $\mathcal{L}_{\mathrm{tc}}$ penalizes Gaussian total correlation, $\mathcal{L}_{\mathrm{hsic}}$ discourages dependence across latent factors, $\mathcal{L}_{\mathrm{aff}}$ is the optional affect classification loss, and $\mathcal{L}_{\mathrm{src}}$ is the optional source-adversarial loss. The multiplicative factors $g_{\bullet}(t) \in [0,1]$ denote the epoch-dependent warmup scales implemented through `reg_scales`.

On top of the base objective, stage-3 training introduces two delayed auxiliary terms:
%
% code map: dcas/pipelines.py:239-257, 258-276, 307-363
\begin{align}
\mathcal{L}_{\mathrm{total}}(t)
=
\mathcal{L}_{\mathrm{base}}
 + \lambda_{\mathrm{pair}}\, \gamma_{\mathrm{pair}}(t)\, \mathcal{L}_{\mathrm{pair}}
 + \lambda_{\mathrm{rank}}\, \gamma_{\mathrm{rank}}(t)\, \mathcal{L}_{\mathrm{rank}},
\end{align}
%
where $\gamma_{\mathrm{pair}}(t)$ and $\gamma_{\mathrm{rank}}(t)$ are stage activation functions determined by `constraint_start_epoch`, `constraint_warmup_epochs`, `rank_start_epoch`, and `rank_warmup_epochs`.
The pairwise term refines $\mathbf{z}^a$ using pseudo-PAL or human PAL constraints, while the ranking term uses interaction-derived rank examples with sampled negatives.

This hierarchy matters for faithful writing. In the implementation, reconstruction and KL terms are the primary optimization backbone, disentanglement and invariance terms are regularizers, and pairwise/ranking terms are late-stage auxiliary supervision rather than co-equal losses from epoch one. Accordingly, the method should be described as a curriculum-based approximate optimization procedure trained with AdamW, not as a globally optimal joint solver.

For the current `V4 main + CultureMERT` stage-3 setting, the nonzero high-level weights are
%
% code map: configs/train/train_v4_main_culturemert_stage3.run.json:11-31
\begin{align}
(\lambda_{\mathrm{pair}}, \lambda_{\mathrm{rank}}, \lambda_{\mathrm{dom}}, \lambda_{\mathrm{con}}, \lambda_{\mathrm{cov}}, \lambda_{\mathrm{tc}}, \lambda_{\mathrm{hsic}}, \lambda_{\mathrm{src}}, \beta_{\mathrm{KL}})
=
(0.15, 0.12, 0.50, 0.20, 0.05, 0.05, 0.02, 0.10, 1.0).
\end{align}
%
We report these values here because they are explicit design choices rather than framework defaults, and because later ablation and sensitivity sections interpret their effect on the quality-fairness trade-off.

